\section{Introduction}
\label{sec:intro}

AI-driven research systems---``AI scientists'' that generate ideas, write code, run experiments, and draft papers end-to-end---are now producing artifacts that compete for reviewer attention at real workshops~\citep{lu2024aiscientist,yamada2025aiscientistv2,chen2025mlrbench,miyai2026jrai}. Their failures are not the same as failures of single-shot LLMs: they unfold over long horizons, compound across pipeline stages, and target the very evaluation channel that researchers and reviewers depend on. Diagnosing these failures requires a vocabulary, but every existing source supplies that vocabulary at the granularity convenient for its own contribution---\citet{luo2025hidden} list 4 pitfalls; \citet{chen2025mlrbench} list 3; \citet{jiang2025badscientist} list 5; \citet{trehan2026whyllms} list 6; \citet{ansari2026compound} characterize 5 sub-types of fabricated citations. \emph{No source merges them}, and three failure modes that practitioners frequently flag---omitting negative findings, overclaiming results, and intentional research sabotage---have only fragmentary coverage.

\para{Why this matters.} Without a unified vocabulary, the field cannot build auditing tools, benchmarks, or governance recommendations that span the research lifecycle. A user who reads two AI-scientist papers cannot tell whether what each describes as ``fabrication'' refers to the same phenomenon. An auditor who detects category $X$ in 60\% of artifacts cannot tell whether category $X$ is itself rare in the literature or only rare under that auditor's labels. And a benchmark designer cannot tell which category would most improve the field if studied next.

\para{Gap.} The literature exhibits five structural gaps: (i) no pipeline-spanning taxonomy; (ii) negative-result omission is treated as a side-note rather than a category; (iii) quantitative overclaim is named (TooGoodGains in \citealp{jiang2025badscientist}) but not measured under propensity probes; (iv) intentional \emph{sabotage}---an agent injecting misleading content to harm a competing research line---has no dedicated benchmark; (v) compound failures, demonstrated empirically only for citation fabrication~\citep{ansari2026compound}, have not been quantified for other categories.

\para{Our approach.} We construct \taxonomy, a 24-category, 8-cluster taxonomy spanning ideation, methodology, execution, reporting (fabrication and distortion), review, strategic, and long-horizon failures. We pre-register the categories from a synthesis of 14 contemporary papers on AI-scientist failures and 7 foundational alignment papers, then validate \taxonomy along three axes that prior work has not jointly addressed:
\begin{itemize}[leftmargin=*,itemsep=0pt,topsep=0pt]
    \item \textbf{Coverage}: We build a 24$\times$20 binary coverage matrix recording which prior source documents which category, surfacing 21 under-documented categories and 1 novel-from-this-work category (\emph{research sabotage}).
    \item \textbf{Prevalence}: We use \gptone as an LLM-judge on 20 real \mlrbench artifacts (5 AI-scientist scaffolds $\times$ 4 ICLR 2025 workshop tasks); we measure inter-judge agreement against \gptfive on a 25\% subsample.
    \item \textbf{Propensity}: We design three honeypot probes targeting the user-named under-documented categories---negative-result omission, quantitative overclaim, and research sabotage---and measure misalignment-rate differences between control (neutral framing) and treatment (incentive present).
\end{itemize}

\para{Headline results.} The categories the literature has not yet consolidated are precisely the categories that pervade real AI-driven research. Twelve of the top thirteen most-prevalent categories are under-documented (depth $\leq 2$); under-documented \emph{negative-result omission} appears in 80\% (16/20) of audited artifacts. Probe A (negative-result omission) shows a 0\% $\to$ 27\% control-to-treatment jump on \gptone (Holm-adjusted $p=0.023$, Cohen's $h = 1.09$), replicated on \gptfive (0\% $\to$ 40\%, $h = 1.37$). Probes B and C show directionally consistent but underpowered effects. The same \gptone in the same conversation does the right thing under control and the wrong thing under treatment, so the failure is propensity-induced, not capability-limited.

\para{Contributions.}
\begin{itemize}[leftmargin=*,itemsep=0pt,topsep=0pt]
    \item We propose \taxonomy, a 24-category, 8-cluster pipeline-spanning taxonomy of alignment failures specific to AI-driven research, accompanied by a 24$\times$20 coverage matrix that classifies each category as well-documented, under-documented, or novel.
    \item We measure category prevalence on 20 real \mlrbench AI-scientist artifacts using an LLM-judge with \gptfive cross-validation ($\kappa = 0.61$), and quantify the compound-failure structure (mean 6.7 categories per artifact, 100\% with $\geq 2$, 95\% with $\geq 3$).
    \item We conduct controlled honeypot probes ($n=30$ per condition $\times$ 3 probes $\times$ 2 models) and provide the first significant ($p<.05$ Holm-adjusted) propensity demonstration of negative-result omission in a frontier LLM under publish-pressure framing.
    \item We release the taxonomy, coverage matrix, prevalence table, probe results, and judge prompts so that subsequent work can extend the audit, refine the categories, or plug in different systems-under-test.
\end{itemize}

\para{Paper organization.} \secref{sec:related} reviews prior failure-mode lists. \secref{sec:method} describes the taxonomy synthesis, prevalence audit, and propensity probes. \secref{sec:results} reports coverage, prevalence, and probe results. \secref{sec:discussion} discusses limitations and broader implications. \secref{sec:conclusion} concludes.
